\chapter{What is actually measured}

\textbf{Purpose:} State the pipeline from a deposited structure to a recovered edge, leaving a reader able to say what the number is a number about.
\textbf{Scope:} \texttt{representation/\allowbreak{}structure/\allowbreak{}crystal.\allowbreak{}py}, \texttt{measure/\allowbreak{}shift\_\allowbreak{}agreement.\allowbreak{}py}

\section{From a deposited cell to a grid}

A CIF entry gives the cell lengths, the angles, and the fractional coordinates of the sites in one cell. Four steps turn that into something the detector reads.

\begin{enumerate}
\item \textbf{The angles are checked.} Only cells with right angles are used. An oblique cell would need the grid sheared before a shift along an index meant a shift along an axis, and that shear is a second thing able to go wrong between the published number and the recovered one.
\item \textbf{The cell is tiled.} A single cell holds one period and one period cannot be told from noise. Tiling is capped so the longest side stays at 320 voxels, which keeps a large cell from asking for a grid nothing can hold.
\item \textbf{The sites are quantized} onto a grid of 0.25 angstroms, each voxel carrying an element code and zero where nothing sits. Two atoms landing on one voxel become one voxel. That is reported and not worked around: at this resolution they are one voxel.
\item \textbf{Each axis is asked separately} for its period, and all three answers are kept. One structure therefore contributes up to three independent results.
\end{enumerate}

Nothing at any step tells the detector what it is looking at. It receives a three dimensional array of small integers and an axis index.

\section{Scoring a lag, and why not the tallest one}

For each lag along the axis, the grid is compared with a copy of itself shifted by that lag, and the agreement is the share of positions holding the same code. Sweep the lags and the lattice period agrees where nothing else does.

\textbf{Taking the tallest lag is wrong, and it was taken that way first.} A lattice of period $P$ agrees with itself just as well at $2P$ and at $3P$. Which of those is tallest is settled by noise, so the tallest lag reports a harmonic and not the fundamental. Read that way, published cell edges came back at almost exactly twice their value on ten of eighteen axes.

The correction is to score a candidate on itself together with its multiples, and to require at least two members of that family to be inside the swept range before the candidate is scored at all. Exactly two members are used and never more: a family growing as the candidate shrinks lets a short wrong candidate win by holding more members, since it only has to catch one good lag among them. Herzenbergite, molybdite and one more each failed that way before the family was capped, reading short by a factor near two thirds.

\textbf{This is the harmonic trap, and it is not specific to crystals.} The same failure has bitten three modules in this work independently. Scoring against multiples instead of taking the tallest peak is the general fix, and the crystal case is where it is easiest to see, because the right answer is published and a factor of two is unmistakable.

\section{What a failure would have looked like}

Stated before the numbers, since afterwards it is worth nothing.

A recovered edge missing the published one by more than one voxel would mean the detector is not reading the period. A systematic bias in one direction would mean it is reading the tiling and not the crystal, since the tiling is the only structure present that the mineral did not put there. Both are visible in the reported columns without knowing which entry is which.
